\section{Implementation in \pkg{gretl}}\label{sec:BACE_implementation}
In this section, we describe the code of the \pkg{BACE}
package. First, we will characterize our code and the use
of both: the command line and the GUI interfaces, then we
will present some examples of using the \pkg{BACE}.
Finally, we will show the outputs that are returned.

\subsection{The core of the \pkg{BACE} code}
The \pkg{BACE} package consists of 4 public and 14
private functions. The name of each function starts
with the \cmd{BACE\_} prefix. The public functions are:
\begin{enumerate}
	\item \cmd{BACE()}~see~\ref{sec:BACE()}
	\item \cmd{BACE\_GUI()}~see~\ref{sec:BACE_GUI()}
	\item \cmd{BACE\_Print()}~see~\ref{sec:BACE_Print()}
	\item \cmd{BACE\_Plot()}~see~\ref{sec:BACE_Plot()}
\end{enumerate}

\subsubsection{\cmd{BACE()}}\label{sec:BACE()}
This is a core package function which runs and controls
the main loop. As a consequence, it is a prime interface
to the \pkg{BACE} package. The \cmd{BACE()} function
returns a bundle with all information needed to print
and/or plot results anytime. One can print the results
by calling \cmd{BACE\_Print()} or plot graphs by calling
\cmd{BACE\_Plot()}. The signature of \cmd{BACE()}
function is as follows:
\begin{lstlisting}
function bundle BACE (series Y "Dependent variable",
  list X_list "List of indenpendent variables (X)",
  bundle Options[null])	
\end{lstlisting}
The first two options are mandatory, while the last one
is an optional collection of additional arguments. Those
arguments are defined with the following default values:\label{sec:default_option_bundle}
\begin{lstlisting}
scalar Options["Y_AR_p"] = 0
scalar Options["with_const"] = 1
scalar Options["model_prior"] = 2
scalar Options["avg_model_size"] = avg_model_size_tmp
scalar Options["alpha"] = 0.6
scalar Options["Nrank"] = xmin(2^(k+Y_AR_p_tmp-1), 40)
scalar Options["do_joint"] = 0
scalar Options["Npredict"] = 0
scalar Options["Nrep"] = Nrep
scalar Options["burn"] = 10
scalar Options["verbosity"] = 2
scalar Options["show_progress"] = do_show_progress
scalar Options["do_plot"] = do_plot_tmp
string Options["model_type"] = "linear"
scalar Options["data_trans"] = 0
scalar Options["bic_type"] = 1
scalar Options["hyperparameter"] = 0.5 # gamma for EBIC
scalar Options["binary_threshold"] = 0.5
scalar Options["low"] = 0
scalar Options["high"] = NA
scalar Options["_gui"] = 0	
\end{lstlisting}
where \textit{k} stands for the number of independent variables,
\lstinline!scalar avg_model_size! is calculated according to
the following formula:
\begin{equation}
	\texttt{\lstinline!avg\_model\_size!} = \left\{
	\begin{array}{lll}
		\textrm{\lstinline!(k + Y\_AR\_p\_tmp)/2!} & \mathrm{if} & \textrm{\lstinline!with_const \!= 2!},\\
		\textrm{\lstinline!(k + Y\_AR\_p\_tmp + 1)/2!} & \mathrm{elsewhere}, &\\
	\end{array}
	\right.
	\label{eq:avg_model_size_setup}
\end{equation}
\lstinline!scalar Nrep! takes the following values:
\begin{equation}
	\texttt{\lstinline!Nrep!} = \left\{
	\begin{array}{ccl}
		10^4 & \mathrm{if} & k \leq 15,\\
		10^5 & \mathrm{if} & k \in [16;19],\\
		10^6 & \mathrm{if} & k \geq 20,\\
	\end{array}
	\right.
	\label{eq:Nrep_setup}
\end{equation}
\lstinline!scalar show_progress! takes the following values:
\begin{equation}
	\texttt{\lstinline!show\_progress!} = \left\{
	\begin{array}{cll}
		1 & \mathrm{if} & \textrm{\lstinline!$sysinfo!\texttt{\small.gui\_mode}},\\
		0 & \mathrm{elsewhere}, &\\
	\end{array}
	\right.
	\label{eq:show_progress_setup}
\end{equation}
and \lstinline!scalar do_plot! takes the following values:
\begin{equation}
	\texttt{\lstinline!do\_plot!} = \left\{
	\begin{array}{cll}
		1 & \mathrm{if} & \textrm{\lstinline!$sysinfo!\texttt{\small.gui\_mode}},\\
		0 & \mathrm{elsewhere}. &\\
	\end{array}
	\right.
	\label{eq:do_plot_setup}
\end{equation}
Most of the optional parameters listed above are of
numerical type and they may be gathered into following
groups according to what kind of numbers they can take%
\footnote{Note, that the \lstinline!scalar _gui! is
	for internal use only and must not be modified
	by user at any rate.}:
\begin{enumerate}
	\item Integers: \lstinline!Y_AR_p!, \lstinline!with_const!,
	\lstinline!model_prior!, \lstinline!Nrank!, \lstinline!do_joint!,
	\lstinline!Npredict!, \lstinline!Nrep!, \lstinline!burn!,
	\lstinline!verbosity!, \lstinline!data_trans!, \lstinline!bic_type!.
	\item Booleans: \lstinline!show_progress!, \lstinline!do_plot!.
	\item Scalars: \lstinline!avg_model_size!, \lstinline!alpha!,
	\lstinline!binary_threshold!.
\end{enumerate}
Optional parameters control the following package options:
\begin{enumerate}
	\item \textbf{Y\_AR\_p}: autoregression order of variable Y.
	\item \textbf{with\_const}: whether we include constant term
		in each estimated model?; possible values are: 0---without
		constant, 1---constant is always included, 2---constant
		may be dropped.
	\item \textbf{model\_prior}: model prior distribution to use;
		possible values are: 0---none, 1---Binomial, 2---Beta-Binomial.
	\item \textbf{avg\_model\_size}: expected average model size
		for model prior.
	\item \textbf{alpha}: significance level for the initial model.
	\item \textbf{Nrank}: number of top ranked models.
	\item \textbf{do\_joint}: do we perform jointness analysis?;
		possible values are: 0---none, 1---Doppelhofer-Weeks
		Measure (DW1), 2---Doppelhofer-Weeks Measure (DW2),
		3---Ley-Steel Measure (LS1), 4---Ley-Steel Measure (LS2),
		5---Strachan Measure, 6---Yule's Q Measure (YQ),
		7---Modified Yule's Q Measure (YQM).
	\item \textbf{Npredict}: number of out-of-sample forecasts.
		Note, that using this option makes \pkg{BACE} to perform
		an analysis on subsample.
	\item \textbf{Nrep}: total number of Monte Carlo iterations
		(model draws).
	\item \textbf{burn}: percentage of burn-in draws in \MCcube.
	\item \textbf{verbosity}: verbosity of the output: 1---basic
		results, 2---full results.
	\item \textbf{show\_progress}: do we show overall progress?
	\item \textbf{do\_plot}: do we produce graphs?
	\item \textbf{model\_type}: model type we should apply to
		the data; possible options are: 1---linear, 2---logit,
		3---probit, 4---poisson.
	\item \textbf{data\_trans}: shall the \pkg{BACE} transform
		data before analysis?; possible options are: 0---none,
		1---centering, 2---standarizing.
	\item \textbf{bic\_type}: Bayesian information criterion
		version which should be used; possible values are:
		1---bic, 2---hbic, 3---ibic, 4---ebic.
	\item \textbf{hyperparameter}: hyperparameter for prior
		setup; for now it is interpreted as $\gamma$ for ebic.
	\item \textbf{binary\_threshold}: applicable only for
		binary-dependent models; this value indicates the
		lowest probability for which $\hat{y}_i$ should take 1.
	\item \textbf{low}: applicable only for	\lstinline!tobit!
		model; this option controls tobit's \lstinline!--llimit!
		flag for setting left bound in censoring.
	\item \textbf{high}: applicable only for \lstinline!tobit!
		model; this option controls tobit's \lstinline!--rlimit!
		flag for setting right bound in censoring.
\end{enumerate}
Some of the above-mentioned parameters are over-loadable, which
means that their value may be passed either via an integer or
via string (case insensitive). Those parameters (together with
value-string mappings) are:\label{sec:over-loadable_options}
\begin{description}
	\item[\textbf{bic\_type}]: 'bic'=1, 'hbic'=2, 'ibic'=3, 'ebic'=4;
	\item[\textbf{model\_prior}]: 'none'=0, 'binomial'=1,
		'betabinomial'=2 (or 'beta-binomial'=2);
	\item[\textbf{with\_const}]: 'never'=0, 'always'=1,
		'canbedropped'=2 (or 'can be dropped'=2);
	\item[\textbf{do\_joint}]: 'none'=0, 'dw1'=1, 'dw2'=2,
		'ls1'=3, 'ls2'=4, 'sm'=5, 'yq'=6, 'myq'=7;
	\item[\textbf{data\_trans}]: 'none'=0, 'centering'=1,
		'standarizing'=2;
	\item[\textbf{model\_type}]: 'linear'=1, 'logit'=2,
		'probit'=3, 'poisson'=4;
\end{description}


\subsubsection{\cmd{BACE\_GUI()}}\label{sec:BACE_GUI()}
The \cmd{BACE\_GUI()} function is an interface to
the \cmd{BACE()} its signature is as follows:
\begin{lstlisting}
function bundle BACE_GUI (
  int model_type[1:5:1] "Model type" {"logit (including ordered)",
  	"probit (including ordered)", "poisson", "tobit",
  	"linear"},
  series Y "Dependent variable",
  int Y_AR_p[0::0] "Y lags (time-series only)",
  list X_list[null] "List of independent variables (X)",
  int with_const[0:2:1] "Constant in model" {"Never", "Always",
  	"Can be dropped"},
  int model_prior[0:2:1] "Model prior" {"none", "Binomial",
  	"Beta-Binomial"},
  int bic_type[1:4:1] "Bayesian information criterion" {"BIC",
  	"HBIC", "IBIC", "EBIC"},
  int do_joint[0:7:0] "Jointness analysis" {"None",
  	"Doppelhofer-Weeks Measure (DW1)", "Doppelhofer-Weeks
	Measure (DW2)", "Ley-Steel Measure (LS1)", "Ley-Steel
	Measure (LS2)", "Strachan Measure (SM)", "Yule's Q
	Measure (YQ)", "Modified Yule's Q Measure (MYQ)"},
  int Npredict[0::0] "Number of out-of-sample forecasts",
  int verbosity[1:2:2] "Verbosity",
  bool show_progress[1] "Show overall progress",
  bool do_plot[1] "Plot graphs?")
\end{lstlisting}	
Note, that although \cmd{BACE\_GUI()} is designed to be
a wrapper to \cmd{BACE()} it does not allow for setting
all possible parameters. This is partially because some
of them are mutually exclusive, partially because we wanted
to design GUI interface which is not very complicated.


\subsubsection{\cmd{BACE\_Print()}}\label{sec:BACE_Print()}
\cmd{BACE\_Print()} is a function for printing results
of analysis performed by \pkg{BACE} package and its
signature is as follows:
\begin{lstlisting}
function void BACE_Print (bundle *MODEL)
\end{lstlisting}
This function takes as an argument reference to the
bundle \lstinline!MODEL! with \pkg{BACE} results and prints is
according to \texttt{verbosity} settings.


\subsubsection{\cmd{BACE\_Plot()}}\label{sec:BACE_Plot()}
\cmd{BACE\_Plot()} is a function for producing graphs
and its signature is as follows:
\begin{lstlisting}
function void BACE_Plot (const bundle *MODEL)
\end{lstlisting}
This function takes as an argument reference to the MODEL
bundle with \pkg{BACE} results and plots graphs according
to \lstinline!do_plot! settings. In non-GUI mode (when
\lstinline!$sysinfo!\texttt{\small.gui\_mode == 0}), plots
are generated as scripts for \texttt{gnuplot} program and
saved on a disk. Names of generated files have their base
set to \lstinline!strftime($now[1], "%Y%m%d%H%M%S")!. 
